\documentclass[11pt]{article}

\usepackage[utf8]{inputenc}
\usepackage[T1]{fontenc}
\usepackage[french]{babel}
\usepackage{lmodern}
\usepackage{geometry}
\usepackage{setspace}
\usepackage{graphicx}
\usepackage{titlesec}
\usepackage{enumitem}
\usepackage{hyperref}
\usepackage{fancyhdr}
\usepackage{amsmath}
\usepackage{amssymb}
\usepackage{tikz}
\usetikzlibrary{arrows.meta,positioning}

\geometry{margin=2.5cm}
\setlength{\parindent}{0pt}
\setlength{\parskip}{0.4em}
\onehalfspacing

% Numérotation simple en bas de page
\pagestyle{fancy}
\fancyhf{}
\fancyfoot[C]{\thepage}
\renewcommand{\headrulewidth}{0pt}

\begin{document}

\section{Définitions et Propriérés}

Toutes les variables aléatoires de ce travail sont définies sur l'espace probabilisé $(\Omega, \mathcal{F}, \mathbb{P})$.
\subsection{Définition d'une chaîne de Markov et propriétés}

\textbf{Définition d'une chaîne de Markov :}\\
Soit $E$ un ensemble fini ou dénombrable, appelé \textit{espace d'états}. Soit $(X_n)_{n \ge 0}$ un \textit{processus} à valeur dans $E$.
On dit que $(X_n)_{n \ge 0}$ est une \textit{chaîne de Markov} d'ordre 1, si :\\
Pour tout $n \ge 0$, pour tout $y \in E$, pour tout $x_0,...,x_{n-1},x \in E$ tel que $\mathbb{P}(X_0 = x_0, \ldots, X_{n-1} = x_{n-1}, X_n = x) > 0$, on a:
\[
\mathbb{P}(X_{n+1} = y \mid X_n = x_n, X_{n-1} = x_{n-1},\ldots,X_0=x_0)=\mathbb{P}(X_{n+1}=x_{n+1} \mid X_n=x_n)
\]
On pose $\pi_0$ la distribution initiale, i.e., la loi de $X_0$.

De plus, s'il existe $P : E \times E \to [0,1]$ tel que $(x,y) \mapsto P(x,y)=\mathbb{P}(X_{n+1} = y \mid X_n = x_n)$, pour tout $n \ge 0$, pour tout $x,y \in E$ tel que $\mathbb{P}(X_n=x_n) > 0$.
On dit alors que $(X_n)_{n \ge 0}$ est une chaîne de Markov d'ordre 1 \textit{homogène}.
Dans ce cas, $P$ est appelé \textit{matrice de transition} du processus $(X_n)_{n \ge 0}$, avec $P(x,y)>0$ et $\sum_{y \in E}P(x,y)=1$, pour tout $x,y \in E$.\\

Dans ce travail, nous ne considèrerons qu'uniquement les chaînes de Markov d'ordre 1 homogènes, c'est pourquoi nous les appelerons juste chaînes de Markov.

\subsection{Propriétés}
\textbf{$\bullet$ Équation de Chapman-Kolmogorov :} \\
Soit $(X_n)_{n \ge 0}$ une chaîne de Markov de matrice de transition P et de distribution initiale $\pi_0$. On a,\\
Pour tout $n \ge 0$, pour tout $x_0, \ldots, x_n \in E$, 
\[
\mathbb{P}(X_0=x_0, \ldots, X_n=x_n) = \pi_0(x_0) \mathbb{P}(x_0,x_1) \ldots \mathbb{P}(x_{n-1},x_n)
\]
i.e. la loi d'une chaîne de Markov est entièrement caractérisée par sa distribution initiale et sa matrice de transition.

\textbf{$\bullet$ Propriété de Markov :} \\
Soit $n \ge 0$, $\pi_0$ une mesure de probabilité sur $E$ et $x \in E$ tel que $\mathbb{P}_{\pi_0}(X_n=x)>0$.\\
Alors, pour tout $A \in E^n$, $B \in E^m$, $m \ge 0$, on a :
\[
\begin{aligned}
\mathbb{P}_{\pi_0}\bigl((X_0,\ldots,X_{n-1}) \in A,\,(X_{n+1},\ldots,X_{n+m}) \in B \mid X_n=x\bigr)
&= \mathbb{P}_{\pi_0}\bigl((X_0,\ldots,X_{n-1}) \in A \mid X_n=x\bigr) \\
&\qquad \times \mathbb{P}_{\pi_0}\bigl((X_{n+1},\ldots,X_{n+m}) \in B \mid X_n=x\bigr).
\end{aligned}
\]
i.e. le passé et le futur d'un moment $n$ d'une chaîne de Markov sont indépendants entre eux.

\newpage

\textbf{$\bullet$ Quelques propriétés utiles :} \\


\end{document}
